\chapter{System Design}
\label{ch:design}

This chapter presents the architecture that satisfies the requirements of
\cref{ch:requirements}. \Cref{sec:architecture-overview} gives a high-level,
layered view of the system. \Cref{sec:mobile-architecture} details the
mobile application's internal architecture. \Cref{sec:local-storage-design}
designs the local database schema. \Cref{sec:backend-design} describes the
backend structure. \Cref{sec:sync-design} is the core of the chapter: it
designs the synchronization algorithm, its triggers, and the
last-write-wins conflict-resolution policy. \Cref{sec:connectivity-design}
describes connectivity monitoring, and \cref{sec:design-traceability}
traces every design element back to the requirement(s) it satisfies.

\section{Architectural Overview}
\label{sec:architecture-overview}

CampusFix follows a layered architecture (\cref{fig:layered-architecture})
with four layers: a \emph{presentation layer} (Compose Multiplatform UI and
ViewModels), a \emph{domain/business layer} embedded in the ViewModels and
the synchronization module for this application's modest scope, a
\emph{data layer} (Room-backed local database), and a
\emph{synchronization layer} that bridges the data layer to the remote
backend. This mirrors the standard mobile app architecture recommended for
Kotlin Multiplatform and Android (a pragmatic variant of Clean
Architecture~\cite{martin2017clean, androidarch}), simplified because the
application's domain logic is thin enough that a separate use-case layer
would add indirection without adding value.

\begin{figure}[htbp]
  \centering
  \begin{tikzpicture}[
      layer/.style={draw, rounded corners, minimum width=10.5cm, minimum height=1.3cm, align=center, font=\small},
      >=Latex
    ]
    \node[layer, fill=blue!8] (pres) at (0,4.5) {\textbf{Presentation layer}\\ Compose Multiplatform screens (Login, Report list, Add report, Report detail) + ViewModels};
    \node[layer, fill=green!8] (data) at (0,3.0) {\textbf{Data layer}\\ Room database (\texttt{ReportDao}, \texttt{UserDao}) over SQLite (via \texttt{sqlite-bundled})};
    \node[layer, fill=violet!8] (sync) at (0,1.5) {\textbf{Synchronization layer}\\ \texttt{SyncManager} + \texttt{ConnectivityObserver}};
    \node[layer, fill=orange!10] (backend) at (0,0) {\textbf{Backend}\\ Firebase Authentication + Firestore (+ Storage for photos)};

    \draw[<->, thick] (pres) -- (data);
    \draw[<->, thick] (data) -- (sync);
    \draw[<->, thick] (sync) -- (backend);
  \end{tikzpicture}
  \caption{Layered architecture of CampusFix. The presentation layer only
  ever talks to the data layer; it has no direct dependency on the
  network or on Firebase, which keeps every screen usable offline by
  construction.}
  \label{fig:layered-architecture}
\end{figure}

The defining architectural rule, and the one from which most of the
offline-first properties follow, is: \textbf{the presentation layer depends
only on the data layer.} No screen or ViewModel calls Firestore directly;
every read comes from Room, and every write goes to Room. The
synchronization layer depends on both the data layer and the backend, and
is the only component that talks to Firestore. This means that removing
the network (or Firestore being unreachable) can never make the UI throw
an exception, hang, or show stale data it does not know is stale --
because the UI's only data source, Room, is always available and always
correct with respect to what has actually been saved on-device.

\section{Mobile Application Architecture}
\label{sec:mobile-architecture}

\subsection{Module and package structure}

The mobile client is a single Kotlin Multiplatform module
(\texttt{composeApp}) with three source sets: \texttt{commonMain}, which
contains all UI, ViewModel, data, and synchronization code shared between
platforms; \texttt{androidMain}, which contains the small amount of
Android-specific glue (activity entry point, platform \texttt{Context}
plumbing, Android's \texttt{ConnectivityManager}-based network callback);
and \texttt{iosMain}, which contains the analogous iOS glue (view
controller entry point, \texttt{NWPathMonitor}-based network callback).
\Cref{tab:package-structure} summarizes the package layout inside
\texttt{commonMain}.

\begin{table}[htbp]
\centering
\caption{Package structure of \texttt{commonMain}.}
\label{tab:package-structure}
\begin{tabularx}{\textwidth}{@{}l X@{}}
\toprule
\textbf{Package} & \textbf{Responsibility} \\
\midrule
\texttt{data.model} & Data classes shared by the local and remote
representations (\texttt{Report}, \texttt{User}, \texttt{SyncStatus}). \\
\texttt{data.local} & Room database, DAOs, type converters, and the local
session store. \\
\texttt{sync} & \texttt{SyncManager} and the \texttt{ConnectivityObserver}
interface (with platform-specific \texttt{actual} implementations). \\
\texttt{viewmodel} & One ViewModel per screen, owning UI state and
orchestrating calls to the data and sync layers. \\
\texttt{ui.screens} & Compose UI for each screen (login, sign-up, report
list, add report, report detail). \\
\texttt{ui.navigation} & Declarative route definitions for
Navigation-Compose. \\
\texttt{util} & Cross-cutting helpers: dispatcher provider (for
platform-appropriate coroutine dispatchers), network reachability check,
password hashing. \\
\bottomrule
\end{tabularx}
\end{table}

\subsection{Presentation pattern: MVVM with unidirectional data flow}

Each screen follows the Model-View-ViewModel (MVVM) pattern, standard for
both Jetpack Compose and Compose Multiplatform applications. A ViewModel
exposes UI state as a Kotlin \texttt{StateFlow}, computed from a
\texttt{Flow} of data read from Room; the Composable screen collects that
state and renders it, and forwards user events (button taps, text input)
back to the ViewModel as method calls. \Cref{fig:mvvm-flow} shows this for
the report list screen, which is representative of the pattern used
throughout the application.

\begin{figure}[htbp]
  \centering
  \begin{tikzpicture}[
      node distance=1.6cm,
      box/.style={draw, rounded corners, minimum width=3.0cm, minimum height=1cm, align=center, font=\small},
      >=Latex
    ]
    \node[box, fill=blue!8] (screen) {ReportListScreen\\(Composable)};
    \node[box, fill=violet!10, right=of screen] (vm) {MainViewModel};
    \node[box, fill=green!8, right=of vm] (dao) {ReportDao.\\getByOwnerFlow()};
    \node[box, fill=orange!10, below=of dao] (db) {Room / SQLite};

    \draw[->] (screen.north) to[out=60,in=120] node[above,font=\scriptsize]{user events} (vm.north);
    \draw[->] (vm.south) to[out=240,in=300] node[below,font=\scriptsize]{StateFlow<UiState>} (screen.south);
    \draw[->] (vm) -- node[above,font=\scriptsize]{Flow<List<Report>>} (dao);
    \draw[<->] (dao) -- (db);
  \end{tikzpicture}
  \caption{Unidirectional data flow for the report list screen. Room's
  \texttt{Flow}-returning queries automatically re-emit whenever the
  underlying table changes (e.g. after a sync updates a row's status),
  which is what makes the UI reactively reflect sync progress without any
  manual refresh logic.}
  \label{fig:mvvm-flow}
\end{figure}

\subsection{Dependency wiring}

Rather than a full dependency-injection framework -- unnecessary at this
project's scale -- CampusFix uses a single hand-written composition root,
\texttt{AppContainer}, constructed once per process from the
platform-supplied \texttt{AppDatabase} instance. \texttt{AppContainer}
owns the DAOs, the \texttt{SyncManager}, the \texttt{ConnectivityObserver},
and the session store, and hands them down to ViewModels. This guarantees
that every screen shares the exact same \texttt{SyncManager} and
\texttt{ConnectivityObserver} instances, which matters because it is the
single \texttt{ConnectivityObserver} subscription that drives automatic
synchronization application-wide (\cref{sec:sync-design}); duplicating it
per-screen would duplicate sync attempts.

\section{Local Storage Design}
\label{sec:local-storage-design}

\subsection{Choice of local database}

The project brief allows SQLite via either a direct binding (e.g. SQLDelight
or the platform SQLite driver) or via Room, Android's higher-level
persistence library. CampusFix uses \textbf{Room over SQLite}, with the
Kotlin Multiplatform-compatible \texttt{room-runtime} artifact and the
bundled cross-platform SQLite driver (\texttt{sqlite-bundled}), which
allows the same Room entities, DAOs, and generated implementation to run
unmodified on both Android and iOS. Room was preferred over a
lower-level driver for three reasons: (1) compile-time verified SQL via
annotation processing (KSP) catches schema/query mismatches before
runtime; (2) DAO methods can return Kotlin \texttt{Flow}s that
automatically re-query and emit on table changes, which is exactly the
reactive behaviour \cref{fig:mvvm-flow} relies on; (3) Room's schema
export (\texttt{room \{ schemaDirectory(...) \}}) gives a versioned,
inspectable record of the database schema for free.

\subsection{Entity-relationship design}

\Cref{fig:er-diagram} shows the local schema. It intentionally contains
fields that have no equivalent on the server (\texttt{id},
\texttt{sync\_status}) and omits fields that only exist during upload
(there is no separate outbox/queue table: a \emph{pending report is simply
a row of the \texttt{reports} table whose \texttt{sync\_status} is
\texttt{PENDING}}). This design choice -- using one table with a status
column rather than a separate queue table -- was made deliberately: it
keeps the "list of reports" and "list of pending work" views trivially
consistent with each other (they are the same table filtered differently),
at the cost of mixing two concerns (application data and sync bookkeeping)
in one entity. \Cref{sec:design-traceability} revisits this trade-off.

\begin{figure}[htbp]
  \centering
  \begin{tikzpicture}[
      entity/.style={draw, minimum width=6.6cm, align=left, font=\small},
      title/.style={font=\small\bfseries},
    ]
    \node[entity] (report) {
      \begin{tabular}{@{}ll@{}}
        \multicolumn{2}{l}{\textbf{reports}} \\
        \midrule
        \underline{id} & INTEGER (PK, autoincrement) \\
        local\_id & TEXT (UUID) \\
        remote\_id & TEXT, nullable \\
        title & TEXT \\
        category & TEXT \\
        description & TEXT \\
        location & TEXT \\
        photo\_path & TEXT, nullable \\
        remote\_photo\_url & TEXT, nullable \\
        owner & TEXT (FK $\rightarrow$ users.email) \\
        created\_at & INTEGER (epoch ms) \\
        updated\_at & INTEGER (epoch ms) \\
        sync\_status & TEXT (enum) \\
      \end{tabular}
    };
    \node[entity, below=1.4cm of report] (user) {
      \begin{tabular}{@{}ll@{}}
        \multicolumn{2}{l}{\textbf{users}} \\
        \midrule
        \underline{email} & TEXT (PK) \\
        password\_hash & TEXT \\
        salt & TEXT \\
        display\_name & TEXT, nullable \\
      \end{tabular}
    };
    \draw[-{Latex[open]}] (report.south west) -- ++(-0.6,0) |- (user.west)
      node[midway,above,font=\scriptsize,sloped]{owner (many-to-one)};
  \end{tikzpicture}
  \caption{Entity-relationship diagram of the local SQLite schema, as
  materialized by Room from the \texttt{Report} and \texttt{User} data
  classes.}
  \label{fig:er-diagram}
\end{figure}

\subsection{Data access layer}

Access to the \texttt{reports} table is mediated entirely through
\texttt{ReportDao}, an interface annotated with Room's \texttt{@Dao}. Two
query shapes are used throughout the app: suspend functions that return a
single snapshot (used for one-off operations, e.g. \enquote{get all
pending reports to sync now}) and \texttt{Flow}-returning functions (used
to drive UI state, e.g. \enquote{observe this user's reports as they
change}). \Cref{tab:dao-methods} lists the DAO surface.

\begin{table}[htbp]
\centering
\caption{Methods exposed by \texttt{ReportDao}.}
\label{tab:dao-methods}
\begin{tabularx}{\textwidth}{@{}l X@{}}
\toprule
\textbf{Method} & \textbf{Purpose} \\
\midrule
\texttt{insert(report)} & Insert a new report; replaces on primary-key
conflict (used when re-saving after obtaining a \texttt{remote\_id}). \\
\texttt{update(report)} & Update an existing row (used after edits and
after synchronization changes its status). \\
\texttt{delete(report)} & Hard-delete a row (used once a
\texttt{deleted\_pending} report's deletion has been confirmed remotely). \\
\texttt{getByOwnerFlow(owner)} & Reactive list of a user's reports,
consumed by the report list screen. \\
\texttt{getByIdFlow(id)} & Reactive single-report stream, consumed by the
report detail screen. \\
\texttt{getPending()} & Snapshot of all rows with
\texttt{sync\_status = PENDING}; the sync module's upload queue. \\
\texttt{getDeletedPending()} & Snapshot of all rows with
\texttt{sync\_status = DELETED\_PENDING}; the sync module's deletion queue. \\
\bottomrule
\end{tabularx}
\end{table}

\section{Backend Structure}
\label{sec:backend-design}

The backend is Firebase, used as a Backend-as-a-Service rather than a
custom API server, per the project brief. Two Firebase products are used:

\begin{itemize}
  \item \textbf{Firebase Authentication} manages user accounts (email and
  password) and issues the ID token used to identify the calling user to
  Firestore. The mobile client caches the authentication state locally
  (\cref{sec:offline-req}, OFR-2) so that a previously-authenticated user
  is not forced back to the login screen when offline.
  \item \textbf{Cloud Firestore} stores one document per report in a
  top-level \texttt{reports} collection, plus a \texttt{users} collection
  for profile metadata not already covered by Firebase Authentication. A
  document's ID is, by design, the same UUID as the report's local
  \texttt{local\_id} whenever possible (\cref{sec:sync-design}), which
  keeps the mapping between local rows and remote documents a pure
  identity function rather than requiring a separate lookup table.
\end{itemize}

\Cref{tab:firestore-doc} shows the shape of a Firestore report document.
It matches the local schema (\cref{fig:er-diagram}) closely by design, with
two differences: the purely local bookkeeping fields \texttt{id} (the
Room auto-increment key) and \texttt{sync\_status} are not uploaded, since
they have no meaning on the server -- \texttt{sync\_status} in particular
is a property of the \emph{local copy's relationship to the server}, not a
property the server itself needs to track.

\begin{table}[htbp]
\centering
\caption{Shape of a Firestore \texttt{reports/\{docId\}} document.}
\label{tab:firestore-doc}
\begin{tabularx}{\textwidth}{@{}l X@{}}
\toprule
\textbf{Field} & \textbf{Type} \\
\midrule
localId & string \\
title & string \\
description & string \\
category & string \\
location & string \\
photoPath & string or null \\
remotePhotoUrl & string or null \\
owner & string \\
createdAt & number (epoch ms) \\
updatedAt & number (epoch ms) \\
\bottomrule
\end{tabularx}
\end{table}

Firestore's own offline persistence cache (enabled by default in most
Firebase SDKs) is deliberately \emph{not} relied upon as the
application's offline story: CampusFix uses Room as the single source of
truth precisely so that offline behaviour does not depend on an
implementation detail of the Firestore SDK that the application does not
control. Firestore is treated purely as a remote synchronization target.

\section{Synchronization Strategy Design}
\label{sec:sync-design}

This is the central design problem of the thesis. \Cref{sec:sync-statuses}
defines the state machine every report moves through.
\Cref{sec:sync-triggers} defines when synchronization runs.
\Cref{sec:sync-algorithm} defines the algorithm itself.
\Cref{sec:conflict-design} defines the conflict-resolution policy.

\subsection{Sync status state machine}
\label{sec:sync-statuses}

Every report's \texttt{sync\_status} field follows the state machine in
\cref{fig:sync-state-machine}.

\begin{figure}[htbp]
  \centering
  \begin{tikzpicture}[
      state/.style={draw, rounded corners, minimum width=2.6cm, minimum height=0.9cm, align=center, font=\small},
      >=Latex
    ]
    \node[state, fill=yellow!20] (pending) at (0,0) {PENDING};
    \node[state, fill=green!15] (synced) at (5,0) {SYNCED};
    \node[state, fill=red!15] (failed) at (0,-2.4) {FAILED};
    \node[state, fill=orange!15] (delpending) at (5,-2.4) {DELETED\_PENDING};
    \node[draw=none] (gone) at (9.2,-2.4) {\emph{(row deleted)}};

    \draw[->] (pending) to[out=20,in=160] node[above,font=\scriptsize]{upload OK} (synced);
    \draw[->] (pending) to[out=-100,in=100] node[left,font=\scriptsize]{upload error} (failed);
    \draw[->] (failed) to[out=80,in=-160] node[below,font=\scriptsize]{retry, OK} (synced);
    \draw[->] (failed) to[out=0,in=180] node[above,font=\scriptsize,align=center]{retry,\\error} (failed.north west) ;
    \draw[->] (pending) -- node[left,font=\scriptsize]{user deletes} (delpending);
    \draw[->] (synced) -- node[right,font=\scriptsize]{user deletes} (delpending);
    \draw[->] (delpending) -- node[above,font=\scriptsize]{remote delete OK} (gone);
  \end{tikzpicture}
  \caption{State machine of the \texttt{sync\_status} field. A report is
  created directly into \texttt{PENDING} and is removed from the local
  database only after a confirmed remote deletion, never as a side effect
  of a failed upload.}
  \label{fig:sync-state-machine}
\end{figure}

The state machine encodes OFR-4 and OFR-8 directly: there is no transition
out of \texttt{PENDING} or \texttt{FAILED} other than a successful or
failed sync attempt, and no transition deletes a row except a confirmed
remote deletion -- a network error can move a report \emph{into}
\texttt{FAILED} but can never destroy it.

\subsection{Synchronization triggers}
\label{sec:sync-triggers}

Synchronization is triggered by three, non-exclusive sources, satisfying
OFR-5 and FR-9:

\begin{enumerate}
  \item \textbf{Connectivity transition.} The \texttt{ConnectivityObserver}
  exposes a \texttt{Flow<Boolean>} of online/offline state; the
  application collects this flow at the top level (in the composition
  root/app entry point) and calls \texttt{SyncManager.sync()} whenever the
  value transitions from \texttt{false} to \texttt{true}.
  \item \textbf{App foreground / cold start.} On process start, after the
  database and \texttt{AppContainer} are initialized, a sync attempt is
  made opportunistically -- it will simply no-op quickly if there is no
  connectivity or nothing pending.
  \item \textbf{Manual user action.} The user can pull-to-refresh the
  report list or tap a retry affordance on a \texttt{failed} report, both
  of which call \texttt{SyncManager.sync()} directly.
  \end{enumerate}

A background periodic sync (e.g. via \texttt{WorkManager} on Android) is
identified as a natural extension in \cref{ch:conclusions} but is not part
of the current design: at this application's usage pattern (short,
foreground sessions to file a report), connectivity-triggered and
app-open-triggered sync were judged sufficient to meet NFR-1--NFR-5, and
adding a periodic background job would introduce additional platform-
specific complexity (battery-optimization exemptions, background task
scheduling APIs differ substantially between Android and iOS) for limited
benefit at this scale.

\subsection{Synchronization algorithm}
\label{sec:sync-algorithm}

\Cref{fig:sync-sequence} shows the sequence of a synchronization run
triggered by a connectivity transition. \Cref{alg:sync} states the same
logic as pseudocode.

\begin{figure}[htbp]
  \centering
  \begin{tikzpicture}[
      participant/.style={font=\small\bfseries},
      lifeline/.style={dashed},
      msg/.style={-{Latex}, font=\scriptsize},
    ]
    \foreach \name/\x in {Connectivity/0, SyncManager/3, ReportDao/6, Firestore/9} {
      \node[participant] at (\x,6) {\name};
      \draw[lifeline] (\x,5.7) -- (\x,0);
    }
    \draw[msg] (0,5.2) -- node[above]{online = true} (3,5.2);
    \draw[msg] (3,4.7) -- node[above]{getPending()} (6,4.7);
    \draw[msg] (6,4.2) -- node[above]{List<Report>} (3,4.2);
    \draw[msg] (3,3.5) -- node[above,align=center]{for each report:\\ set(doc)} (9,3.5);
    \draw[msg] (9,3.0) -- node[above]{ack / error} (3,3.0);
    \draw[msg] (3,2.4) -- node[above,align=center]{update(report,\\ status=SYNCED/FAILED)} (6,2.4);
    \draw[msg] (3,1.7) -- node[above]{getDeletedPending()} (6,1.7);
    \draw[msg] (6,1.2) -- node[above]{delete(doc) / delete(row)} (9,1.2);
  \end{tikzpicture}
  \caption{Sequence of a synchronization run: pending reports are read from
  Room, pushed to Firestore one by one, and their local status updated
  according to the outcome; deletions pending remote confirmation are
  processed afterwards.}
  \label{fig:sync-sequence}
\end{figure}

\begin{figure}[htbp]
\begin{center}
\begin{minipage}{0.92\textwidth}
\begin{quote}
\ttfamily\small
\textbf{function} sync(): \\
\hspace*{1em} pending $\leftarrow$ dao.getPending() \\
\hspace*{1em} uploaded $\leftarrow$ 0 \\
\hspace*{1em} \textbf{for} report \textbf{in} pending: \\
\hspace*{2em} docId $\leftarrow$ report.remoteId \textbf{or} report.localId \\
\hspace*{2em} \textbf{try}: \\
\hspace*{3em} firestore.collection("reports").document(docId).set(report) \\
\hspace*{3em} dao.update(report.copy(remoteId = docId, status = SYNCED, updatedAt = now())) \\
\hspace*{3em} uploaded $\leftarrow$ uploaded + 1 \\
\hspace*{2em} \textbf{catch} error: \\
\hspace*{3em} \textit{leave report as PENDING/FAILED; log; continue with next report} \\
\hspace*{1em} \textbf{for} report \textbf{in} dao.getDeletedPending(): \\
\hspace*{2em} \textbf{try}: \\
\hspace*{3em} firestore.collection("reports").document(report.remoteId \textbf{or} report.localId).delete() \\
\hspace*{3em} dao.delete(report) \\
\hspace*{2em} \textbf{catch} error: \textit{leave as DELETED\_PENDING; continue} \\
\hspace*{1em} \textbf{return} uploaded
\end{quote}
\end{minipage}
\end{center}
\caption{Synchronization algorithm (pseudocode). Using the report's own
\texttt{local\_id} as the Firestore document ID whenever no
\texttt{remote\_id} exists yet makes \texttt{set()} idempotent: re-running
the same upload after an interrupted sync overwrites the same document
rather than creating a duplicate, which is how OFR-6 (resumability) and
NFR-7 (idempotency) are satisfied without an explicit checkpoint or
two-phase commit.}
\label{alg:sync}
\end{figure}

Two design decisions in this algorithm deserve explicit justification:

\begin{description}
  \item[Per-report try/catch, not an all-or-nothing batch.] Wrapping each
  report's upload in its own error handler means one malformed or
  oversized report cannot block every other pending report from
  synchronizing. This is important at scale: with a queue of 50--100
  reports (the \enquote{large queue} scenario of
  \cref{ch:evaluation}), a single failure must degrade gracefully rather
  than stall the entire batch.
  \item[Client-chosen document ID.] Letting the client derive the
  Firestore document ID from \texttt{local\_id} (rather than letting
  Firestore auto-generate one on first write) is what makes the upload
  operation naturally idempotent, satisfying OFR-6 with no additional
  bookkeeping. The trade-off is a very small, application-scale increase
  in the probability of a document-ID collision, which is negligible
  given UUIDv4 identifiers.
\end{description}

\subsection{Conflict resolution design}
\label{sec:conflict-design}

Per FR-12/OFR-7 and the project brief's guidance, CampusFix adopts
\textbf{last-write-wins (LWW) based on \texttt{updated\_at}} as its
conflict-resolution policy: when a report exists in two versions, the
version with the greater \texttt{updated\_at} timestamp is retained and
the other discarded. LWW is a well-known, simple point in the space of
conflict-resolution strategies (as opposed to, e.g., automatic
field-level merging, operational transformation, or CRDTs, all of which
are strictly more powerful but also more complex to implement and to
reason about~\cite{shapiro2011crdt,terry1995bayou}); it is an
appropriate choice here specifically because reports are
\textbf{single-owner} records (\cref{sec:constraints}) with low edit
frequency, so the practical cases in which LWW silently discards a change
a user cares about are rare, and the user can always re-edit.

\Cref{fig:conflict-scenario} illustrates the target scenario the policy is
designed for: the same report, identified by \texttt{local\_id}, edited on
two devices while one of them was offline.

\begin{figure}[htbp]
  \centering
  \begin{tikzpicture}[
      ev/.style={draw, rounded corners, minimum width=3.0cm, minimum height=0.8cm, align=center, font=\scriptsize},
      >=Latex
    ]
    \node[ev, fill=blue!8] (a1) at (0,2) {Device A (online)\\edits title, $t=10$};
    \node[ev, fill=orange!10] (b1) at (0,0) {Device B (offline)\\edits description, $t=12$};
    \node[ev, fill=green!10] (server) at (5.5,1) {Firestore doc\\updatedAt=10 after A};
    \node[ev, fill=violet!10] (b2) at (5.5,-1.4) {Device B reconnects,\\ pushes its version, $t=12$};
    \node[ev, fill=green!15] (final) at (10.4,-0.2) {Firestore doc\\updatedAt=12\\(B's version wins)};

    \draw[->] (a1) -- (server);
    \draw[->] (b1) -- (b2);
    \draw[->] (server.east) to[out=-20,in=160] (b2.west);
    \draw[->] (b2) -- (final);
  \end{tikzpicture}
  \caption{Last-write-wins conflict scenario: Device A's online edit
  (title, $t{=}10$) reaches the server first, but Device B's offline edit
  (description, $t{=}12$) has a later timestamp and overwrites it entirely
  once B reconnects -- including the title change, which is lost. This is
  the accepted trade-off of field-level-unaware LWW, discussed further in
  \cref{sec:conflict-eval}.}
  \label{fig:conflict-scenario}
\end{figure}

As \cref{fig:conflict-scenario} makes explicit, whole-record LWW loses A's
title change even though the two edits touched different fields and could,
in principle, both have been kept. A field-level merge would avoid this,
at the cost of needing to track per-field modification timestamps and a
merge step with no natural place to run (Firestore, used purely as a
document store here, does not execute application logic server-side
without adding Cloud Functions, which is out of scope). Whole-record LWW
was chosen as the simplest policy that is \emph{unambiguous} and
\emph{terminating} -- properties that matter more than optimality for a
first implementation -- and the trade-off is documented here and revisited
against real conflict-test results in \cref{sec:conflict-eval}.

It is important to be precise about what \emph{is} and \emph{is not}
currently implemented (this distinction is discussed further as a
limitation in \cref{ch:conclusions}): the design above specifies a
symmetric, bidirectional LWW merge between local and remote versions. The
current implementation of \texttt{SyncManager} (\cref{ch:implementation})
realizes the \emph{local-to-remote push} half of this design --
uploading pending local changes and always overwriting the remote
document with \texttt{set()} -- but does not yet \emph{pull} remote
changes back down and compare \texttt{updated\_at} before overwriting.
Practically, this means the current build behaves as \enquote{last local
write wins on push,} which is sufficient for the offline-creation and
delayed-sync scenarios of \cref{ch:evaluation}, but does not yet fully
realize OFR-7 for the two-device edit-conflict scenario; closing this gap
is scoped as the highest-priority future-work item in
\cref{sec:future-work}.

\section{Connectivity Monitoring Design}
\label{sec:connectivity-design}

Connectivity is modeled behind a single \texttt{ConnectivityObserver}
interface (\cref{sec:mobile-architecture}) exposing
\texttt{val isOnline: Flow<Boolean>}, with platform-specific
implementations: Android's implementation wraps
\texttt{ConnectivityManager.NetworkCallback}, and iOS's wraps
\texttt{NWPathMonitor} from the Network framework. Both implementations
are required to emit the current state immediately upon collection and
thereafter emit only on actual changes (i.e. the flow is expected to be
naturally deduplicated by the underlying OS callback semantics), so that a
\texttt{SyncManager} collecting the flow reacts exactly once per
offline$\to$online transition rather than repeatedly. This expect/actual
pattern is the only place in the entire synchronization path that requires
platform-specific code; \texttt{SyncManager} itself is 100\% common code.

\section{Traceability}
\label{sec:design-traceability}

\Cref{tab:traceability} maps each offline-specific requirement of
\cref{sec:offline-req} to the design element(s) that satisfy it, closing
the loop between \cref{ch:requirements} and this chapter.

\begin{table}[htbp]
\centering
\caption{Traceability from offline-specific requirements to design.}
\label{tab:traceability}
\begin{tabularx}{\textwidth}{@{}lX@{}}
\toprule
\textbf{Requirement} & \textbf{Design element} \\
\midrule
OFR-1 (local-first writes) & Layered architecture (\cref{fig:layered-architecture}): presentation depends only on the data layer. \\
OFR-2 (offline session) & Local \texttt{SessionStore}, consulted before any network auth call. \\
OFR-3 (deterministic local identity) & \texttt{local\_id} UUID, generated client-side at creation (\cref{fig:er-diagram}). \\
OFR-4 (explicit sync state) & \texttt{sync\_status} state machine (\cref{fig:sync-state-machine}). \\
OFR-5 (connectivity-triggered sync) & \texttt{ConnectivityObserver} + trigger in \cref{sec:sync-triggers}. \\
OFR-6 (resumable sync) & Client-chosen, idempotent document IDs (\cref{alg:sync}). \\
OFR-7 (conflict policy) & LWW policy, \cref{sec:conflict-design} (bidirectional merge: designed, push-only: implemented). \\
OFR-8 (no data loss on failure) & Per-report error handling leaves rows \texttt{PENDING}/\texttt{FAILED}, never deletes on error (\cref{fig:sync-state-machine}). \\
\bottomrule
\end{tabularx}
\end{table}
